% ДҮГНЭЛТ
\chapter*{ДҮГНЭЛТ}
\addcontentsline{toc}{chapter}{Дүгнэлт}
\label{Conclusion}
\pagecolor{white}

Энэхүү дипломын ажилд QR кодоор суурилсан рестораны ширээний мобайл захиалгын
системийг загварчилж, хэрэгжүүлж, туршилтаар баталгаажуулав. Судалгааны зорилт
бүрийн дагуу хүрсэн үр дүн:

\textbf{З1 --- Онолын үндэслэл:} QR код, WebSocket, RESTful API, JWT болон
React Native технологиудын онолыг нарийвчлан судлаж, олон улсын 30 гаруй
бүтээлийн дүн шинжилгээгээр санал болгосон системийн онолын суурийг тогтоов.
Kim нар \cite{kim2023}, Zhang \& Liu \cite{zhang2022}, Patel нар \cite{patel2022}
зэрэг судлаачдын дүгнэлтүүдтэй нийцсэн онолын загварчлал боловсруулав.

\textbf{З2 --- Шаардлагын тодорхойлолт:} Ресторан эздэнтэй болон зочидтой
хийсэн ярилцлагад тулгуурлан 6 функционал болон 4 функционал бус шаардлагыг
тодорхойлов. Хэрэглэгчийн персон загварчлал нь дизайны шийдлүүдэд удирдамж болов.

\textbf{З3 --- Системийн загварчлал:} 4 давхаргат архитектур, ER диаграм, класс
диаграм, дарааллын диаграм, бизнес процессын диаграм, дэлгэцийн зохиомж болон
RESTful API-ийн бүрэн тодорхойлолтыг боловсруулав. Загварчлалд аюулгүй байдлын
RBAC хандалтын удирдлага, OWASP зарчмуудыг харгалзав.

\textbf{З4 --- Хэрэгжилт:} React Native + Expo мобайл клиент, Node.js/Express API
сервер, PostgreSQL өгөгдлийн сан, Socket.IO WebSocket серверийн уялдаа бүхий бүрэн
системийг хэрэгжүүллээ. Захиалга үүсгэснээс ажилтны дэлгэц мэдэгдэл авах хүртэлх
дундаж хугацаа 8 секунд байлаа.

\textbf{З5 --- Туршилт ба харьцуулалт:} Улаанбаатар хотын 3 рестораны 1,240
захиалгад дүн шинжилгээ хийж, дараах гол үр дүнг олж авав:

\begin{itemize}
    \item Захиалгын нийт хугацаа \textbf{80.8\%} (308 → 59 секунд) буурсан;
    \item Захиалгын алдааны түвшин \textbf{91.4\%} (23.3\% → 2.0\%) буурсан;
    \item Ажилтны мэдэгдэл авах дундаж хугацаа \textbf{42 мс} (WebSocket);
    \item Нийт хэрэглэгчийн сэтгэл ханамж \textbf{4.5/5.0};
    \item Зочдын \textbf{91.7\%} тус системийг ашигладаг ресторанд дахин ирнэ.
\end{itemize}

\medskip

Эдгээр үр дүн нь QR дээр суурилсан мобайл захиалгын систем нь Монгол рестораны
орчинд практикт нэвтрүүлэх боломжтой, хэмжигдэхүйн чадамжтай шийдэл болохыг
баталгаажуулж байна. Монгол хэлний дэмжлэг, QPay болон SocialPay нийцтэй байдал,
хямд нэвтрүүлэлтийн зардал нь гадаадын системүүдийн нөхцөн гүйцэтгэж чадаагүй
хэрэгцээг хангасан.

\medskip

\textbf{Ирээдүйн чиглэл:}
\begin{enumerate}
    \item \textbf{Offline горим:} Сүлжээ тасарсан үед захиалгыг локал хадгалж,
          дахин холбогдоход автоматаар дамжуулах;
    \item \textbf{AI зөвлөмж:} Захиалгын түүхт суурилсан хоолны санал болгох;
    \item \textbf{Олон рестораны платформ:} SaaS загварт шилжүүлэх;
    \item \textbf{Хот/хөдөөгийн харьцуулсан судалгаа:} Монгол улсын орон нутгийн
          ресторанд нэвтрүүлэх нөхцөлийн судалгаа.
\end{enumerate}
